\section{One scoring core, two retention protocols}\label{sec:design}

\subsection{The object scored}\label{sec:object}

In both systems the predictand is the quantity that management reads, not a fitted index of it: assessed SSB in kilotonnes for cod, and the calendar-year mean of daily-high J-17 elevation in feet above mean sea level for the aquifer. Each candidate forecaster is a \emph{module} on a \emph{ladder}: a forward-ordered set of models in which each rung adds one structural element to the rung below, so that the contribution of each element can be isolated by comparing a module with its comparator (the next-simpler module) and not only with a naive benchmark. Two baselines are scored on every origin set: last-value persistence, $\hat y_{t+h}=y_t$, and the training mean, $\hat y_{t+h}=\bar y_{1:t}$ (``climatology''). Baselines are always computed on exactly the origin set used for the modules at that horizon, so that every comparison is origin-matched.

\subsection{Scores, origins and horizons}\label{sec:scores}

The primary score is the root mean squared error of the point forecast in the predictand's own units; mean absolute error is reported alongside, and indicator (Brier) scores for the management threshold are secondary and play no part in retention (SI S7, S19). Two evaluation designs are used in both systems. \emph{Fixed windows} fit once on a training window and forecast the following test window as a path; they are diagnostic, and show how a module behaves across an event (a collapse, a recovery, a drought of record) whole. The \emph{rolling-origin} design \citep{Tashman2000} is the one the retention protocols are applied to: at every origin $t$ from the first admissible year to the last that leaves $h$ target years, each module is refitted on the data through $t$ and a forecast of $y_{t+h}$ is scored; the errors are pooled across origins. Horizons are $h=1$ and $h=5$ years in both systems (the horizons at which annual advice and multi-year planning are issued); the horizon profile $h=1,\dots,5$ of the two baselines is added in Section~\ref{sec:synthesis}. Minimum training lengths are 8 years (cod, Specification A), 12 years (cod, Specification B) and 15 years (J-17), giving origin counts of 25/21, 59/55 and 75/71 at $h=1$/$h=5$ (SI S2, S16).

\subsection{The information set at the origin}\label{sec:info}

A forecast issued at origin $t$ for target $t+h$ may use only quantities dated $t$ or earlier, with declared exceptions that are the structural claim under test. On cod, the stock-flow modules receive the realised catch $C_{t+1},\dots,C_{t+h}$ along the horizon: the forecast is conditional on removals, the form in which surplus-production projections are issued (``if removals are $C$, biomass will be $S$''), and it gives those modules information the baselines never see, so that the persistence benchmark is conservative for them. The autonomous cod modules hold catch at its training mean. On J-17 the causal modules must substitute something for the unknown year-$(t+1)$ fluxes: last year's values (M2), the training-mean fluxes (M2m), or a one-step recharge forecast from climate indices known in year $t$ (climate-informed variants); an \emph{oracle} that receives the realised fluxes is scored as a diagnostic and excluded from retention. The contrast between the two systems is deliberate. On cod the driver (catch) is supplied; on J-17 the driver (recharge) must be forecast. Section~\ref{sec:synthesis} shows that this difference decides which term of the error decomposition each system fails on.

\subsection{The two retention protocols}\label{sec:protocols}

Each system is judged by the protocol under which it was scored. Both protocols share a scoring core, coded before the first scoring pass of the respective system and applied unchanged: compare the module's rolling-origin RMSE with origin-matched persistence and with its comparator, on declared windows and origin sets. They differ in the decision predicate.

\begin{center}
\fbox{\parbox{0.95\textwidth}{\small
\textbf{Protocol C (cod).} A module $M$ with comparator $M'$ is \emph{retained} on a specification if and only if, in the rolling-origin design,
\begin{enumerate}[label=(C\arabic*),leftmargin=3em,itemsep=1pt]
\item $\RMSE(M) < 0.95\,\RMSE(M')$ at both $h=1$ and $h=5$ (comparator hurdle; vacuous for the first rung), and
\item $\RMSE(M) < 0.95\,\RMSE(\pers)$ at both $h=1$ and $h=5$ (persistence hurdle), and
\item both hold on the same specification; a win on Specification A does not carry to B or conversely.
\end{enumerate}
The 5\% band treats differences within it as ties.

\medskip
\textbf{Protocol E (Edwards).} A causal module $M$ is \emph{retained} if and only if, at $h=1$ on the point rolling-origin RMSE,
\begin{enumerate}[label=(E\arabic*),leftmargin=3em,itemsep=1pt]
\item $\RMSE(M) < \RMSE(\pers)$, and
\item $\RMSE(M) < \RMSE(M_{\text{next-simpler causal}})$.
\end{enumerate}
No tie tolerance and no minimum margin is imposed; diagnostic oracles and the service series are excluded from retention.}}
\end{center}

Both protocols are decision rules on point estimates; the uncertainty layer of Section~\ref{sec:uncertainty} describes their margins but is not part of either. Protocol C is the stricter, and when the J-17 scores are restated under it no module is retained (Section~\ref{sec:edw_verdict}); Protocol E retains the aquifer's AR(1). We report both, because the purpose is not to make the two systems agree but to show what each rule does with the same scores.

\paragraph{What was completed after scores were seen.} Neither protocol was prospectively registered, and the honest description is a fixed computational protocol with disclosed completions. On cod, the 5\% band and the comparator declarations for the non-nested modules were formalised after the first pass had been scored; neither decides any outcome, the smallest pooled deficit being 17.1\% (SI S1 records the pass order). On J-17, the frozen protocol documents (dated 2026-08-25) contain the retention predicate verbatim, but three elements sit outside it and are recorded as deviations: a \emph{class clause} by which the climatological-flux map M2m, which satisfies (E1) and (E2), is listed and then declined because under constant fluxes it reduces to the AR(1) function class; the declaration of the retained AR(1), rather than the declined M2m, as the comparator for the climate-informed rung; and a struck direction-of-change score for which no value is reported (SI S15). The class clause is the one completion with a visible consequence, and Section~\ref{sec:edw_verdict} reports the verdict with and without it.

\paragraph{No pooling.} The two systems are never pooled, concatenated or averaged; nor are the two cod specifications, whose 1983--2015 overlap differs by 126\,kt RMSE between vintages. Every number in this paper belongs to one series under one protocol, and the cross-system statements of Section~\ref{sec:synthesis} are statements about mechanisms, not about pooled skill.

\subsection{Uncertainty of the margins}\label{sec:uncertainty}

For every load-bearing module--benchmark pair the per-origin squared-error differences $d_t = e^2_{M,t}-e^2_{B,t}$ are summarised by the Diebold--Mariano statistic \citep{Diebold1995} with an autocorrelation-consistent variance truncated at lag $h-1$, the form appropriate to $h$-step errors, reported as a $z$-score without the small-sample correction of \citet{Harvey1997}. It is not the primary uncertainty statement: samples are small, the loss differences are heavy-tailed because a few event-year errors dominate the sums of squares, and the asymptotic reference overstates significance in that setting \citep{Diebold2015}. The primary statement is a paired moving-block bootstrap \citep{Kunsch1989} of the RMSE gap, with the same block indices applied to both loss series; block length $\max(h,3)$ and 20\,000 replicates on cod, block length 8 and 10\,000 replicates on J-17 with an independent second implementation at block length $h$ and 20\,000 replicates (SI S3, S18). Block-length sensitivity is reported in both systems and no conclusion depends on it.

\subsection{What the test could have seen}\label{sec:power}

A non-retention can mean that the structure is absent or that the test could not see it, and the two systems supply the distinction by different means. On cod, seven data-generating processes were simulated at the fitted parameters of the fixed windows (200 replicates each, two process-noise levels): the autonomous truth at collapse- and recovery-window parameters, the stock-flow truth, the depensatory truth, a persistence-true null, and two truths outside the class; each replicate was scored by Protocol C, with pre-declared bars of power $\ge 0.80$ and false retention $\le 0.10$ (SI S8). On J-17 the corresponding evidence is the oracle: the same one-pool map, given the realised fluxes of the target year, scores 7.55\,ft against persistence's 13.23\,ft, so the mechanism is demonstrably able to explain the increment when its driver is known, and the forecasting failure is located in the driver rather than in the map (Section~\ref{sec:edw_mech}). The two designs answer the same question, whether the instrument is capable of a positive result, and in both systems the answer is yes.
